\subsection{Topology on the set of formal power series. Summable families}

By definition A{[}{[}I{]}{]} is nothing other than the product set \(\mathbf{A}^{\mathbf{N}^{(I)}}\) . Except for express mention to the contrary we shall equip A with the discrete topology and A{[}{[}I{]}{]} with the product topology (Gen.~Top. I, p.~31 f.) which we shall call the canonical topology. Equipped with addition and the discrete topology, A is a separated and complete topological group; hence for addition A{[}{[}I{]}{]} is a separated and complete topological group (Gen.~Top., III, p.~238 and 242 and Gen.~Top., II, p.~187). Moreover the algebra A{[}(X \(_{i}\) ) \(_{i\in I}\) {]} of polynomials is dense in A{[}{[}I{]}{]} (Gen.~Top., III p.~238, Prop. 25) and we may thus consider A{[}{[}I{]}{]} as the completion of A{[}(X \(_{i}\) ) \(_{i\in I}\) {]}.

For each \(\beta \in \mathbf{N}^{(I)}\) let \(S_{\beta}\) be the set of multi-indices \(\nu\) such that \(\nu \leqslant \beta\) and let \(\alpha_{\beta}\) be the set of formal power series \(u = \sum_{\nu} \alpha_{\nu} X^{\nu}\) such that \(\alpha_{\nu} = 0\) for

\(\nu \in S_{\beta}\) . Clearly \(S_{\beta}\) is a finite subset of \(\mathbf{N}^{(I)}\) , and every finite subset of \(\mathbf{N}^{(I)}\) is contained in a set of the form \(S_{\beta}\) . It follows that the family \((\mathfrak{a}_{\beta})_{\beta \in \mathbf{N}^{(I)}}\) is a fundamental system of neighbourhoods of 0 in A{[}{[}I{]}{]}. The sets \(\mathfrak{a}_{\beta}\) are ideals in A{[}{[}I{]}{]}, hence (Gen.~Top., III, p.~275) A{[}{[}I{]}{]} is a topological ring.

Lemma 1. --- Let L be an infinite set and \((u_{\lambda})_{\lambda \in L}\) a family of elements of A{[}{[}I{]}{]}, and put \(u_{\lambda} = \sum_{\nu} \alpha_{\lambda, \nu} X^{\nu}\) for \(\lambda \in L\) . Then the following conditions are

equivalent :

\begin{enumerate}
\def\labelenumi{(\roman{enumi})}
\item
  The family \((u_{\lambda})_{\lambda \in L}\) is summable (Gen.~Top., III, p.~262) in A{[}{[}I{]}{]}.
\item
  We have \(\lim u_{\lambda} = 0\) , taken along the filter of complements of finite subsets of \(L\) .
\item
  For every \(\nu \in \mathbf{N}^{(I)}\) we have \(\alpha_{\lambda, \nu} = 0\) except for a finite number of indices \(\lambda \in L\) .
\end{enumerate}

When these conditions hold, the series \(u = \sum_{\lambda \in L} u_{\lambda}\) is equal to \(\sum_{\nu} \alpha_{\nu} X^{\nu}\) with \(\alpha_{\nu} = \sum_{\lambda \in L} \alpha_{\lambda, \nu}\) for each \(\nu \in \mathbf{N}^{(I)}\) .

The equivalence of (i) and (ii) follows from Cor. 2 of Gen.~Top., III, p.~263. The equivalence of (ii) and (iii) follows from the properties of limits in a product space (Gen.~Top., I, p.~55, Cor. 1).

The last assertion follows from Prop. 4 of Gen.~Top., III, p.~266.

Let us give some examples of summable families.

\begin{enumerate}
\def\labelenumi{\alph{enumi})}
\item
  Let \(u \in \mathbf{A}[[\mathrm{I}]]\) and let \(\alpha_{\nu}\) be the coefficient of \(\mathbf{X}^{\nu}\) in \(u\) . The family \((\alpha_{\nu}\mathbf{X}^{\nu})_{\nu \in \mathbb{N}^{(I)}}\) is then summable, with sum \(u\) (which justifies writing \(u = \sum_{\nu} \alpha_{\nu}\mathbf{X}^{\nu}\) ).
\item
  Let \(u \in \mathbf{A}[[\mathbf{I}]]\) ; for every integer \(p \geqslant 0\) let \(u_p\) be the homogeneous component of degree \(p\) of \(u\) . Then the family \((u_p)_{p \in \mathbb{N}}\) is summable and we have \(u = \sum_{p \geqslant 0} u_p\) .
\item
  Let \((u_{\lambda})_{\lambda \in L}\) be a family of elements of A{[}{[}I{]}{]} and suppose that for every integer \(n \geqslant 0\) the set of \(\lambda \in L\) such that \(\omega(u_{\lambda}) < n\) is finite. Then the family \((u_{\lambda})_{\lambda \in L}\) is summable.
\end{enumerate}

Remark. --- Suppose that I is finite. For every integer \(n \geqslant 0\) let \(\mathbf{b}_n\) be the set of formal power series \(u \in \mathbf{A}[[\mathbf{I}]]\) such that \(\omega(u) \geqslant n\) . The sequence \((\mathbf{b}_n)_{n \geqslant 0}\) is a fundamental system of neighbourhoods of 0 in A{[}{[}I{]}{]}. Therefore a family of elements \(u_{\lambda}\) of A{[}{[}I{]}{]} ( \(\lambda \in L\) ) is summable if and only if for every \(n \in \mathbb{N}\) the set of \(\lambda \in L\) such that \(\omega(u_{\lambda}) < n\) is finite.

PROPOSITION 1. --- Let \((u_{\lambda})_{\lambda \in L}\) and \((v_{\mu})_{\mu \in M}\) be two summable families of elements of A{[}{[}I{]}{]}. Then the family \((u_{\lambda}v_{\mu})_{(\lambda, \mu) \in L \times M}\) is summable and we have

\[
\sum_ {(\lambda , \mu) \in L \times M} u _ {\lambda} v _ {\mu} = \left(\sum_ {\lambda \in L} u _ {\lambda}\right) \left(\sum_ {\mu \in M} v _ {\mu}\right).\tag{2}
\]

Let \((\alpha_{\lambda, \nu})_{\nu \in \mathbf{N}^{(I)}}\) (resp. \((\beta_{\mu, \nu})_{\nu \in \mathbf{N}^{(I)}}\) ) be the family of coefficients of \(u_{\lambda}\) (resp. \(v_{\mu}\) ). For each \(\nu \in \mathbf{N}^{(I)}\) there exists only a finite number of pairs \((\nu_1, \nu_2) \in \mathbf{N}^{(I)} \times \mathbf{N}^{(I)}\) such that \(\nu_1 + \nu_2 = \nu\) , hence only a finite number of pairs \((\lambda, \mu) \in L \times M\) such that the coefficient of \(X^{\nu}\) in \(u_{\lambda}v_{\mu}\) is \(\neq 0\) . Hence the family \((u_{\lambda}v_{\mu})_{(\lambda, \mu) \in L \times M}\) is summable. Now the formula (2) follows from the associativity of the sum (Gen.~Top., III, p.~265, formula (2)).

In A{[}{[}I{]}{]} the product is an associative and commutative composition law. We may therefore speak of a multipliable family of elements of A{[}{[}I{]}{]} and of the product of a multipliable family (Gen.~Top., III, p.~262, remark 3).

PROPOSITION 2. --- Let \((u_{\lambda})_{\lambda \in L}\) be a summable family of elements of A{[}{[}I{]}{]}. (i) The family \((1 + u_{\lambda})_{\lambda \in L}\) is multipliable.

\begin{enumerate}
\def\labelenumi{(\roman{enumi})}
\setcounter{enumi}{1}
\tightlist
\item
  Let \(\mathfrak{F}\) be the set of all finite subsets of L. For any \(\mathbf{M} \in \mathfrak{F}\) put \(u_{\mathbf{M}} = \prod_{\lambda \in \mathbf{M}} u_{\lambda}\) . Then the family \((u_{\mathbf{M}})_{\mathbf{M} \in \mathfrak{F}}\) is summable and we have
\end{enumerate}

\[
\sum_ {M \in \mathfrak {F}} u _ {M} = \prod_ {\lambda \in L} (1 + u _ {\lambda}).
\]

Let us define the ideals \(\mathfrak{a}_{\beta}\) as at the beginning of this No., and let \(\beta \in \mathbf{N}^{(I)}\) . There exists a finite subset \(L_0\) of \(L\) such that \(u_{\lambda} \in \mathfrak{a}_{\beta}\) for \(\lambda \notin L_0\) . Then for every \(M \in \mathfrak{F}\) such that \(M \subset L_0\) we have \(u_M \in \mathfrak{a}_{\beta}\) . It follows that the family \((u_M)_{M \in \mathfrak{F}}\) is summable. On the other hand, for any finite subset \(M_0\) of \(L\) we have

\[
\sum_ {M \subset M _ {0}} u _ {M} = \prod_ {\lambda \in M _ {0}} (1 + u _ {\lambda}).
\]

Taken along the filtered ordered set \(\mathfrak{F}\) , the left-hand side has as limit \(\sum_{M\in\mathfrak{F}}u_{M}\) . Hence the right-hand side has as limit \(\sum_{M\in\mathfrak{F}}u_{M}\) , which proves (i) and (ii) at

the same time.

PROPOSITION 3. --- Let \(u = \sum_{\nu} \alpha_{\nu} X^{\nu} \in \mathbf{A}[[\mathbf{I}]]\) and \(m\) an integer \(> 0\) . For every \(n \in \mathbb{N}\) let \((\alpha_{\nu, n})_{\nu \in \mathbb{N}^{(I)}}\) be the family of coefficients of \(u^n\) . If \(\alpha_0^m = 0\) , then \(\alpha_{\nu, n} = 0\) for \(n \geqslant |\nu| + m\) .

Let \(\nu \in \mathbf{N}^{(I)}\) and \(n\in \mathbf{N}\) . We have

\[
\alpha_ {\nu , n} = \sum_ {\nu (1) + \dots + \nu (n) = \nu} \alpha_ {\nu (1)} \dots \alpha_ {\nu (n)}.
\]

If \(n \geqslant |\nu| + m\) and \(\nu(1) + \cdots + \nu(n) = \nu\) , we have \(|\nu(1)| + \cdots + |\nu(n)| \leqslant n - m\) . We thus have \(\nu(r) = 0\) and so \(\alpha_{\nu(r)} = \alpha_0\) for at least m distinct values of r; it follows that \(\alpha_{\nu(1)} \ldots \alpha_{\nu(n)} = 0\) , whence the result.

COROLLARY. --- Let \(u \in \mathbf{A}[[\mathbf{I}]]\) ; then for \(\lim_{n \to \infty} u^n = 0\) to hold it is necessary and

sufficient that the constant term of \(u\) should be nilpotent.

Let \(\alpha_0\) be the constant term of \(u\) . The constant term of \(u^n\) is \(\alpha_0^n\) , hence the stated condition is necessary; it is sufficient by Prop. 3.
% semantic-fragments: legacy-input-seam
\relax{}
